\documentclass[10pt]{beamer}

\usepackage[english]{babel}
\usepackage[utf8]{inputenc}
\usetheme{metropolis}

\usepackage{appendixnumberbeamer}
\usepackage{pgfplots}
\usepackage{booktabs}
\usepackage[scale=2]{ccicons}
\usepackage{minted}

\usepackage{xspace}
\newcommand{\themename}{\textbf{\textsc{metropolis}}\xspace}

\usemintedstyle{emacs}

\newcommand\Fontvi{\fontsize{10}{12}\selectfont}

\title{Applicative programming with effects}
% \subtitle{A modern beamer theme}
\date{\today}
\author{
  Luis Eduardo Bueso de Barrio\\
  lbueso@acm.org
}
\graphicspath{{./pics}}
\institute{ACM UPM}
\titlegraphic{
  \hfill\includegraphics[height=1.5cm]{pics/logo_acm_cromo.png}
  % \includegraphics[height=1.5cm]{pics/sig_plig.png}
}

\begin{document}

\maketitle

\begin{frame}{Contents}
  \setbeamertemplate{section in toc}[sections numbered]
  \tableofcontents[hideallsubsections]
\end{frame}

\section{Why would you need Applicative functors?}

\begin{frame}
  When we are dealing with functors sometimes we have this kind of
  structure:

  \inputminted{text}{../src/intro.txt}

  when we find a type like this:

  \inputminted{haskell}{../src/intro_type.hs}
\end{frame}

\section{Introduction}

\begin{frame}{Introduction}
  \begin{quote}
    This is the story of a pattern that popped up time and again in our
    daily work, programming in Haskell, until the temptation to abstract
    it became irresistible. -- Simon Peyton Jones
  \end{quote}
\end{frame}

\begin{frame}{Sequencing commands}
  Sequencing commands and collect the results is a very common
  pattern.

  \inputminted{haskell}{../src/sequence_type.hs}

  Let's take a look at some implementations for this pattern.
\end{frame}

\begin{frame}{Sequencing commands}
  If we implement this sequence operator for the \verb~IO~ we will get
  something like this using do-notation:

  \inputminted{haskell}{../src/sequence.hs}
\end{frame}

\begin{frame}{Sequencing commands}
  We will rewrite it using the \verb~ap~ function from the
  \verb~Monad~ library.

  \inputminted{haskell}{../src/ap.hs}

  So we get something like this:

  \inputminted{haskell}{../src/sequence1.hs}
\end{frame}

\begin{frame}{Transposing 'matrices'}
  If we model matrices as lists of lists, we can implement
  transposition as:

  \inputminted{haskell}{../src/transpose.hs}

  The \verb~repeat~ function is just:

  \inputminted{haskell}{../src/repeat.hs}
\end{frame}

\begin{frame}{Transposing 'matrices'}
  If we implement some \verb~zapp~ function like this:

  \inputminted{haskell}{../src/zapp.hs}

  We can rewrite transpose as:

  \inputminted{haskell}{../src/transpose1.hs}
\end{frame}

\begin{frame}{Comparison between the functions}
  So let's take a look at their types:

  \inputminted{haskell}{../src/comparison_types.hs}

  and their implementation:

  \inputminted{haskell}{../src/comparison_impl.hs}
\end{frame}

\begin{frame}{Comparison between the functions}
  We can see a relation in the ``application'' function:

  \inputminted{haskell}{../src/comparison_types_ap.hs}

  So we can create an abstraction of this operation called ``apply''
  which will be the function:

  \inputminted{haskell}{../src/apply_type.hs}

  All this patterns will lead us to the creation of the
  \verb~Applicative~ type class.
\end{frame}

\section{The Applicative class}

\begin{frame}{The Applicative class}

  \inputminted{haskell}{../src/applicative.hs}

  We can check that applicative functors are functors by implementing
  \verb~fmap~:

  \inputminted{haskell}{../src/fmap.hs}
\end{frame}

\begin{frame}{The Applicative class}

  Any expression can be transformed to ``canonical form''

  \inputminted{text}{../src/canonical.txt}

  We can transform any monad into an applicative just with:

  \inputminted{haskell}{../src/to_monad.hs}
\end{frame}

\begin{frame}{The Applicative class}
  An \verb~Applicative~ instance for \verb~IO~ would be:

  \inputminted{haskell}{../src/io_instance.hs}

  So our \verb~sequence~ function becomes:

  \inputminted{haskell}{../src/sequence_ap.hs}
\end{frame}

\begin{frame}{The Applicative class}
  The same way we have various instances of \verb~Monoid~ for a type,
  we have various of \verb~Applicative~.

  In our list example we need the \verb~Applicative ZipList~ instance:

  \inputminted{haskell}{../src/ListInstance.hs}

  So our \verb~transpose~ function becomes:

  \inputminted{haskell}{../src/transpose_ap.hs}

  It must be notice that \verb~repeat~ are not the \verb~return~ and
  \verb~ap~ equivalents of any \verb~Monad~.

\end{frame}

\begin{frame}{Applicative laws}
  \begin{itemize}
  \item Identity\\
    \verb~pure id <*> u = u~
  \item Composition\\
    \verb~pure (.) <*> u <*> v <*> w = u <*> (v <*> w)~
  \item Homomorphism\\
    \verb~pure f <*> pure x = pure (f x)~
  \item Interchange
    % \verb~u <*> pure x = pure (\f -> f x) <*> u~
  \end{itemize}
\end{frame}

\begin{frame}{QuickCheck}

  \verb~QuickCheck~ can help us writing correct \verb~Applicative~
  functors using \verb~Checkers~

  In order to use it we must implement two instances:

  \inputminted{haskell}{../src/ArbitraryEqClass.hs}

  And to make the test we just need to:

  \inputminted{haskell}{../src/trigger.hs}
\end{frame}

\begin{frame}[standout]
  \Fontvi
  \inputminted{text}{../src/testOut.txt}
\end{frame}

\section{Traversing data structures}

\begin{frame}{Traversing data structures}
  Let's take a look at \verb~sequence~ and \verb~transpose~ types
  again:

  \inputminted{haskell}{../src/trav_typ.hs}

  This common pattern is called \textit{applicative distributor} for
  lists:

  \inputminted{haskell}{../src/dist.hs}

  This is again in \textit{applicative style}.
\end{frame}

\begin{frame}{Traversing data structures}
  This is usually used with \verb~map~, for example here we use it for
  ``failure propagation'':

  \inputminted{haskell}{../src/flakyMap.hs}

  We can implement a generalization of this:

  \inputminted{haskell}{../src/traverse.hs}
\end{frame}

\begin{frame}{Traversing data structures}
  This pattern is very useful so we can abstract in a type class:

  \inputminted{haskell}{../src/traversable.hs}

  If we implement an \verb~Id~ type like:

  \inputminted{haskell}{../src/Id.hs}

  We can implement \verb~fmap~ very easily:

  \inputminted{haskell}{../src/traversable_fmap.hs}

\end{frame}

\begin{frame}{Traversing data structures}
  Another interesting \verb~Traversable~ instance is:

  \inputminted{haskell}{../src/Traversable_tree.hs}

  In the latest \verb~GHC~ we need \verb~Functor~, \verb~Applicative~
  and also \verb~Foldable~ in order to have a \verb~Traversable~
  instance.
\end{frame}

\section{Monoids are phantom Applicative functors}

\begin{frame}{Monoids are phantom Applicative functors}
  As we know the \verb~Monoid~ is just:

  \inputminted{haskell}{../src/Monoid.hs}

  Monoids are very useful in functional programming, we have them in
  numeric types, lists, booleans... And every monoid also induces an
  applicative functor!

\end{frame}

\begin{frame}{Monoids are phantom Applicative functors}
  If we define:

  \inputminted{haskell}{../src/Accy.hs}

  We can implement an applicative functor that accumulates
  computations:

  \inputminted{haskell}{../src/Monoid_applicative.hs}
\end{frame}

\begin{frame}{Monoids are phantom Applicative functors}
  This accumulation can be seen as a special kind of traversal:

  \inputminted{haskell}{../src/acc_red.hs}

  With this monoid instance we get operations like \verb~flatten~ and
  \verb~concat~ nearly for free!

  \inputminted{haskell}{../src/flatten_concat_monoid.hs}
\end{frame}

\section{Applicative versus Monad}

\begin{frame}{Applicative versus Monad}
  We have already seen how we every \verb~Monad~ can be made
  \verb~Applicative~:

  \inputminted{haskell}{../src/to_monad.hs}

  However there exists more \verb~Applicative~ functors than
  \verb~Monads~.

  This is because \verb~Monads~ are more powerful than
  \verb~Applicative~!

  If we take a look at the \verb~bind~ operator:

  \inputminted{haskell}{../src/bind.hs}

  We can see that it can have an effect on the next computation, while
  \verb~<*>~ just can't.

\end{frame}

\begin{frame}{Applicative versus Monad}
  Let's take a look at an example:

  \inputminted{haskell}{../src/miffy_monad.hs}

  \inputminted{haskell}{../src/iffy_ap.hs}

  If we call this functions like this:

  \inputminted{text}{../src/iffy.txt}

  What is going on?

\end{frame}

\begin{frame}{Composing applicative functors}
  \verb~Applicative~ functors are easier to compose than
  \verb~Monads~.

  If we define a data type \verb~:.:~ to compose \verb~Applicative~,
  we get:

  \inputminted{haskell}{"../src/composition.hs"}

  So we can tell that composing two \verb~Applicative~ functors also
  give us an \verb~Applicative~ functor.

  But when we compose two monads we at least get an \verb~Applicative~
  functor.

\end{frame}

\begin{frame}{Accumulating exceptions}
  If we define a data type to model exceptions:

  \inputminted{haskell}{../src/Except.hs}

  If we use a \verb~Monad~ instance, it will abort once the computation
  fails, however we can define the \verb~Applicative~ instance as:

  \inputminted{haskell}{../src/ExceptApplicative.hs}

  With this instance we can collect errors using the list
  \verb~Monoid~
\end{frame}

% \section{Applicative functors, categorically}
\section{Conclusions}

\begin{frame}{Conclusions}
  We have identified \verb~Applicative~ functors. We saw that every
  \verb~Monad~ is also an \verb~Applicative~ functor.

  We also analysed the pattern that will give us the notion of how
  \verb~Applicative~ functors are created, and how to distinguish them
  in terms of types.

  We have seen \verb~Applicative~ functors laws the ~Checkers~ library
  to test our abstractions laws.

  We have made an abstraction over \verb~Applicative~ functors which
  is ~Traversable~ functors.

  We have defined monoids in terms of \verb~Applicative~ functors.

  We have compared them with ~Monads~.
\end{frame}

\appendix

\begin{frame}{References}

  \textit{Applicative programming with effects}
  \href{http://www.staff.city.ac.uk/~ross/papers/Applicative.html}{\beamergotobutton{Link}}

  \textit{Haskell Book}
  \href{http://haskellbook.com/}{\beamergotobutton{Link}}

  \textit{Category Theory for Programmers}
  \href{https://bartoszmilewski.com/2014/10/28/category-theory-for-programmers-the-preface/}{\beamergotobutton{Link}} %

\end{frame}

\begin{frame}[standout]
  Questions?
\end{frame}

\end{document}